%HW02.tex
%
% Second Homework for Honors Algebra I  (415H)
% Frank Sottile
%%%%%%%%%%%%%%%%%%%%%%%%%%%%%%%%%%%%%%%%%%%%%%%%%%%%%%%%%%%%%%%%%%%%%%%
\documentclass[12pt]{article}
\usepackage{multicol,amsfonts, amssymb,  mathtools,amsmath}
\usepackage[dvipsnames]{xcolor}
\usepackage{colordvi,graphicx}
\headheight=8pt
%
\topmargin=-75pt
\textheight=720pt   \textwidth=560pt
\oddsidemargin=-60pt \evensidemargin=-60pt

\pagestyle{empty}   %  Change this for your write-up

%%%%%%%%%%%%%%%%%%%%%%%%%%%%%%%%%%%%%%%%%%%%

\newcommand{\NN}{{\mathbb N}}
\newcommand{\QQ}{{\mathbb Q}}
\newcommand{\RR}{{\mathbb R}}
\newcommand{\ZZ}{{\mathbb Z}}

\newcommand{\vect}[2]{(\begin{smallmatrix}#1\\#2\end{smallmatrix})}

\def\Color#1#2{\special{color push cmyk #1}#2\special{color pop}}
%\def\Indigo#1{\Color{.42 1. 0. .49}{#1}}
\def\Indigo#1{\Color{1. .95 .05 .4}{#1}}
\def\MyViolet#1{\Color{.6 1. 0. .15}{#1}}
\def\TAMU#1{\Color{.15 1. .39 .69}{#1}}

\newcommand{\barsl}{\noindent\begin{minipage}[t]{590pt}
\Indigo{\rule{590pt}{1.2pt}}\vspace{-5.7mm}\\
\MyViolet{\rule{590pt}{1.2pt}}\vspace{-5.7mm}\\
\Blue{\rule{590pt}{1.2pt}}\vspace{-5.7mm}\\
\Green{\rule{590pt}{1.2pt}}\vspace{-5.7mm}\\
\Yellow{\rule{590pt}{1.2pt}}\vspace{-5.7mm}\\
\Orange{\rule{590pt}{1.2pt}}\vspace{-5.7mm}\\
\Red{\rule{590pt}{1.2pt}}\bigskip
\end{minipage}}


\newcommand{\barsn}{\noindent\begin{minipage}[t]{560pt}
\Indigo{\rule{560pt}{1.1pt}}\vspace{-4.5mm}\\
\MyViolet{\rule{560pt}{1.1pt}}\vspace{-4.5mm}\\
\Blue{\rule{560pt}{1.1pt}}\vspace{-4.5mm}\\
\Green{\rule{560pt}{1.1pt}}\vspace{-4.5mm}\\
\Yellow{\rule{560pt}{1.1pt}}\vspace{-4.5mm}\\
\Orange{\rule{560pt}{1.1pt}}\vspace{-4.5mm}\\
\Red{\rule{560pt}{1.1pt}}\bigskip
\end{minipage}}

\def\demph#1{{\color{Maroon}#1}}
\def\defcolor#1{{\color{Maroon}#1}}

\begin{document}
\LARGE 
\noindent
Algebra 415H \ \ Autumn 2026\vspace{2pt}\\
{\color{magenta}Your Name}\vspace{2pt}\\
\Large 27 August 2025 \hfill
\sf
 Second Homework\makebox[40pt][l]{\ }
\normalsize\vspace{10pt}

\noindent
Write your answers neatly, in complete sentences.  
I highly recommend recopying your work before handing it in.
Correct and crisp proofs are greatly appreciated; oftentimes your work can be shortened and made clearer.

\noindent
I also recommend talkng with other students, and I am able to give hints to public posts in Piazza.

\noindent
\barsn

\noindent\Maroon{{\sf Hand in at the start of class, Thursday 3 September:}} 


\begin{enumerate}
\setcounter{enumi}{5}

%%%%%%%%%%%%%%%%%%%%%%%%%%%%%%%%%%%%%%%%%%%%%%%%%%%%%%%%%%%%%%%%%%%%%%%%%%%%%%%%%%%%%%%%%%%%%%%%%%%%
\item   Let $X$ be any nonempty set.
  Define \defcolor{$\mbox{Bij}(X)$} to be the set of bijections $f\colon X\to X$ from $X$ to itself.
  Show that $\mbox{Bij}(X)$ forms a group under function composition.

%%%%%%%%%%%%%%%%%%%%%%%%%%%%%%%%%%%%%%%%%%%%%%%%%%%%%%%%%%%%%%%%%%%%%%%%%%%%%%%%%%%%%%%%%%%%%%%%%%%%

%%%%%%%%%%%%%%%%%%%%%%%%%%%%%%%%%%%%%%%%%%%%%%%%%%%%%%%%%%%%%%%%%%%%%%%%%%%%%%%%%%%%%%%%%%%%%%%%%%%%
\item  {\color{Maroon}\sf Strange multiplicative rationals.}
  Let $n\in\NN^\times$ be a non-zero (and hence positive) natural number.
  Define the binary operation $\defcolor{*_n}\colon \QQ_+\times\QQ_+\to \QQ_+$ on the positive rational numbers by
  \[
  a *_n b\ \vcentcolon=\ nab\,.  \makebox[1.5cm][l]{\quad\qquad(The usual multiplication)}
  \]
  Show that $(\QQ_+, *_n)$ forms a group.
  (For this, you need to find its identity element and prove that every element has an inverse, among other things.)

 
  
%%%%%%%%%%%%%%%%%%%%%%%%%%%%%%%%%%%%%%%%%%%%%%%%%%%%%%%%%%%%%%%%%%%%%%%%%%%%%%%%%%%%%%%%%%%%%%%%%%%%

  
%%%%%%%%%%%%%%%%%%%%%%%%%%%%%%%%%%%%%%%%%%%%%%%%%%%%%%%%%%%%%%%%%%%%%%%%%%%%%%%%%%%%%%%%%%%%%%%%%%%%
\item  Give a simple argument based on group tables to determine (up to isomorphism)
  all groups of order 4 (i.e.~those with four elements).

%%%%%%%%%%%%%%%%%%%%%%%%%%%%%%%%%%%%%%%%%%%%%%%%%%%%%%%%%%%%%%%%%%%%%%%%%%%%%%%%%%%%%%%%%%%%%%%%%%%%

%%%%%%%%%%%%%%%%%%%%%%%%%%%%%%%%%%%%%%%%%%%%%%%%%%%%%%%%%%%%%%%%%%%%%%%%%%%%%%%%%%%%%%%%%%%%%%%%%%%%
\item Let $(G,*,e)$ be a group and $a,b\in G$.
  Show that $(a*b)^{-1}=a^{-1}*b^{-1}$ if and only if $a*b=b*a$.

%%%%%%%%%%%%%%%%%%%%%%%%%%%%%%%%%%%%%%%%%%%%%%%%%%%%%%%%%%%%%%%%%%%%%%%%%%%%%%%%%%%%%%%%%%%%%%%%%%%%

%%%%%%%%%%%%%%%%%%%%%%%%%%%%%%%%%%%%%%%%%%%%%%%%%%%%%%%%%%%%%%%%%%%%%%%%%%%%%%%%%%%%%%%%%%%%%%%%%%%%
\item Suppose that $G$ is a group in which every element is its own inverse.
  (That is, $a\in G\Rightarrow a*a=e$.)
  Show that $G$ is an abelian group in that its group operation $*$ is commutative: for all $a,b\in G$
  we have $a*b=b*a$.

%%%%%%%%%%%%%%%%%%%%%%%%%%%%%%%%%%%%%%%%%%%%%%%%%%%%%%%%%%%%%%%%%%%%%%%%%%%%%%%%%%%%%%%%%%%%%%%%%%%%

%%%%%%%%%%%%%%%%%%%%%%%%%%%%%%%%%%%%%%%%%%%%%%%%%%%%%%%%%%%%%%%%%%%%%%%%%%%%%%%%%%%%%%%%%%%%%%%%%%%%
\item
  Let $n$ be a positive integer and define $\defcolor{B_+}\subset GL(n,\RR)$ to be the set of upper triangular invertible
  matrices.
  (A matrix $M\in B_+$ if $M_{i,j}=0$ whenever $i>j$.)
  Show that $B_+$ forms a group under matrix multiplication.

%%%%%%%%%%%%%%%%%%%%%%%%%%%%%%%%%%%%%%%%%%%%%%%%%%%%%%%%%%%%%%%%%%%%%%%%%%%%%%%%%%%%%%%%%%%%%%%%%%%%

%%%%%%%%%%%%%%%%%%%%%%%%%%%%%%%%%%%%%%%%%%%%%%%%%%%%%%%%%%%%%%%%%%%%%%%%%%%%%%%%%%%%%%%%%%%%%%%%%%%%
\item   Let $\ell_1$ and $\ell_2$ be two lines in the plane $\RR^2$ which meet in a point, $o$.
  Let $\sigma_i\colon \RR^2\to\RR^2$ be the reflection of the plane in the line $\ell_i$.
  Using, say, elementary Euclidean geometry, show that the composition $\sigma_2\circ \sigma_1$ of the reflection in line
  $\ell_1$ followed by the reflection in $\ell_2$  {\color{Magenta}({\bf not} reflection in the image of $\ell_2$)} is rotation through the
  angle $2\theta$ about their point $o$ of intersection, where $\theta$ is the angle from $\ell_1$ to $\ell_2$.
  You can use that these are rigid motions, if that helps.


%%%%%%%%%%%%%%%%%%%%%%%%%%%%%%%%%%%%%%%%%%%%%%%%%%%%%%%%%%%%%%%%%%%%%%%%%%%%%%%%%%%%%%%%%%%%%%%%%%%%
\item  In the previous problem, what is $\sigma_2\circ\sigma_1$ when the lines are parallel?

%%%%%%%%%%%%%%%%%%%%%%%%%%%%%%%%%%%%%%%%%%%%%%%%%%%%%%%%%%%%%%%%%%%%%%%%%%%%%%%%%%%%%%%%%%%%%%%%%%%%


%%%%%%%%%%%%%%%%%%%%%%%%%%%%%%%%%%%%%%%%%%%%%%%%%%%%%%%%%%%%%%%%%%%%%%%%%%%%%%%%%%%%%%%%%%%%%%%%%%%%

%%%%%%%%%%%%%%%%%%%%%%%%%%%%%%%%%%%%%%%%%%%%%%%%%%%%%%%%%%%%%%%%%%%%%%%%%%%%%%%%%%%%%%%%%%%%%%%%%%%%
\item  (Possibly hard question) Is the group $(\QQ_+,*_n)$ isomorphic to, or not isomorphic to, the usual multiplicative group
  $(\QQ_+, \cdot,1)$ of positive rational numbers?
  Justify your answer.

%%%%%%%%%%%%%%%%%%%%%%%%%%%%%%%%%%%%%%%%%%%%%%%%%%%%%%%%%%%%%%%%%%%%%%%%%%%%%%%%%%%%%%%%%%%%%%%%%%%%
    
\end{enumerate}
%%%%%%%%%%%%%%%%%%%%%%%%%%%%%%%%%%%%%%%%%%%%%%%%%%%%%%%%%%%%%%%%%%%%%%%%%%%%%%%%%%%%%%%%%%%%%%%%%%%%

\end{document}

%%%%%%%%%%%%%%%%%%%%%%%%%%%%%%%%%%%%%%%%%%%%%%%%%%%%%%%%%%%%%%%%%%%%%%%%%%%%%%%%%%%%%%%%%%%%%%%%%%%%
\item

%%%%%%%%%%%%%%%%%%%%%%%%%%%%%%%%%%%%%%%%%%%%%%%%%%%%%%%%%%%%%%%%%%%%%%%%%%%%%%%%%%%%%%%%%%%%%%%%%%%%
